% !TEX root = ../../cookbook.tex
\newpage
\recipe{Moussaka}{}{}{}

\begin{multicols}{2}
	\subsection*{Ingredienti}
	
	\subsubsection*{Per il Ragú}
	\begin{ingredients}
		\item 1 kg di carne magra macinata (\sfrac{1}{2} di manzo e \sfrac{1}{2} di maiale)
		\item 100 g di passata di pomodoro
		\item 25 g di concentrato di pomodoro
		\item 1 cipolla
		\item 1 costa di sedano
		\item 2 spicchi d'aglio
		\item \sfrac{1}{2} cucchiaino di zucchero
		\item 1 bicchiere di vino rosso
		\item 2 foglie di alloro
		\item 1 stecca di cannella
		\item Anice stellato q.b.
		\item Prezzemolo q.b.
		\item Noce moscata q.b.
		\item Olio d'oliva q.b.
		\item Sale e pepe q.b.
	\end{ingredients}
	
	\subsubsection*{Per gli strati}
	\begin{ingredients}
		\item 1 kg di melanzane
		\item Se fritte:
		\begin{ingredients}
			\item Farina q.b.
			\item Olio per friggere
		\end{ingredients}
		\item 3 patate (opzionale)
	\end{ingredients}
	
	\subsubsection*{Per la besciamella allo yogurt}
	\begin{ingredients}
		\item 500 g di latte
		\item 250 g di yogurt greco
		\item 50 g di farina
		\item 50 g di burro
		\item Noce moscata q.b.
		\item Pepe q.b.
	\end{ingredients}
	
	\subsubsection*{Per comporre}
	\begin{ingredients}
		\item Parmigiano q.b.
		\item Pangrattato q.b.
		\item Origano secco q.b.
		\item Olio d'oliva q.b.
	\end{ingredients}
	
\end{multicols}

% --- Preparation ---
\textbf{Iniziare col ragú:} \\
Tritare grossolanamente il \textbf{sedano} e la \textbf{cipolla}.
Rosolarli con l'\textbf{aglio} finché appassiscono, poi aggiungere lo \textbf{zucchero}.
Aggiungere la \textbf{carne macinata} e rosolare a fiamma alta, sgranando col mestolo. 
Dopo circa 5 minuti, sfumare col \textbf{vino} e aromatizzae con \textbf{alloro}, \textbf{anice}, e \textbf{cannella}.
Unire la \textbf{passata} e il \textbf{concentrato di pomodoro}, mescolare bene e far sobbollire a fuoco vivace, quindi aggiungere il \textbf{prezzemolo} tritato.
Regolare di \textbf{noce moscata}, \textbf{sale}, e \textbf{pepe}. 
Abbassare la fiamma e cuocere scoperto per $\sim$1.5 h.

\textbf{Passare alle melanzane e patate:} \\ 
Tagliare per il lungo le \textbf{melanzane} a fette di ca. 1 cm, e le \textbf{patate} di ca. \sfrac{1}{2} cm. 
Tradizionalmente, le verdure vanno infarinate e fritte ma, alternativamente, possono essere spennellate con olio d'oliva e cotte in forno per una versione piu semplice e meno oleosa\footnote{In questo caso, cospargere le fette con del sale da entrambi i lati e lasciargli perdere un po' di acqua per ca. \sfrac{1}{2}h, poi asciugarle bene prima di comporre la moussaka.}. 
Tenerle da parte.

\textbf{Fare la besciamella allo yogurt:} \\
Scaldare il \textbf{latte} a parte. 
In un pentolino, sciogliere il burro a fuoco medio-basso e aggiungere la \textbf{farina} per ottenere un roux\footnote{Il \textit{roux} è il connubio cotto di burro e farina in parti uguali tipico della cucina francese.} dorato.
Aggiungere il latte caldo poco a poco mentre si mischia con un frustino per sciogliere eventuali grumi, quindi aromatizzare con \textbf{sale}, \textbf{pepe}, e \textbf{noce moscata}. 
Ottenuta una besciamella corposa, togliere dal fuoco e lasciar raffreddare 5 minuti prima di mischiare con lo \textbf{yogurt}, e infine correggere di sale, pepe, e noce moscata. 

\textbf{Comporre la moussaka:} \\ 
Ungere una pirofila dai bordi alti con dell'\textbf{olio di oliva} (base $\sim$26x20 cm), coprire il fondo con della besciamella, creare una base con tutte le patate, disporre uno strato di melanzane, e spolverizzare con \textbf{parmigiano}, \textbf{pangrattato} e \textbf{origano}.
Ricoprire con tutto il ragú.
Disporre un altro strato di melanzane, spolverizzarle con parmigiano, e pangrattato, e origano.
Coprire tutto con la besciamella, quindi completare con pangrattato e parmigiano.
Cuocere in forno preriscaldato a 180°C per $\sim$45 minuti. 



% --- Description ---
\begin{recipeblock}
	
\end{recipeblock}
